\section{Conclusion}
\label{sec:conclusion}

Reusable natural-language skills give workflow owners an editable point of
control over LLM agents, but the skill text is not the workflow that an agent
will enact. A failed execution identifies an unacceptable outcome without
fully specifying the policy that should replace it, and a plausible textual
revision does not reveal how far its behavioral effects will extend. Skill
repair therefore requires review across user intent, artifact revision, and
runtime behavior.

We introduced \emph{behavioral patch review}, an interaction model that makes
these relationships inspectable before a persistent change is released.
\systemname{} helps an owner construct a provisional behavioral scope across
cases, connects that scope to an evidence-augmented skill diff, and compares
the complete original and revised skills under matched execution. The system
uses AI to support evidence collection and implementation while preserving
unresolved policy and release authority as explicit owner judgments.

Our formative, technical, and controlled-study evidence suggests that this
organization can improve understanding of enacted behavioral change, alignment
with a prospective repair scope, and evidence-grounded release decisions. More
broadly, the work argues that maintaining reusable agents should not be framed
only as prompt editing or automated optimization. The consequential object of
review is the behavioral change a revised artifact may carry into future
workflows, and the people who own that change need representations that make
its intent, implementation, evidence, and limits jointly reviewable.
